\documentclass[11pt]{article}
\usepackage{scriptstyle}

\begin{document}

\begin{center}
{\LARGE\bfseries CollusionGraph --- Presentation Script}\\[4pt]
{\large The Investigator Console}\\[8pt]
{\small Total running time: \textbf{about 21 minutes}\quad\textbar\quad
30 slides\quad\textbar\quad Start here. This is the opening of the talk.}
\end{center}

\vspace{6pt}
\begin{tcolorbox}[colback=metagrey!6,colframe=metagrey!45,boxrule=0.6pt,arc=1pt]
\textbf{How to use this file.} Each slide has three boxes.
\textbf{Say this} is spoken almost word for word --- it is written to sound
like talking, not like reading. \textbf{What this actually means} is for you,
not the room; read it beforehand so you understand what you are saying.
\textbf{If the evaluator asks} is the question that actually gets asked there,
with an answer you can give without leaving the screen.

A ``\bt'' inside a spoken line is a breath. Slow down there.
\do{square brackets} are things you \emph{do} with the mouse, not words you say.

\textbf{Set-up before you start.} Console open on \textbf{Overview}, domain
\textbf{Financial \textbullet\ AML}, dataset \textbf{elliptic\_pp}, review
number \textbf{50}. Copilot closed. Nothing else on screen.
\end{tcolorbox}

% ===========================================================================
\part{Part 1 --- The opening (2 minutes)}

\slide{Slide 1 --- Opening --- Nobody Looks Guilty On Their Own}{HIGH (8/10)}{$\approx$ 55 seconds}

\begin{saythis}
Good morning everyone. We're Group \blank.

I want to start with the one thing that makes this project hard.

Pick any single bank account inside a money-laundering ring. Just one. On its
own, it looks completely ordinary. Money comes in, money goes out, and it
hasn't broken a single rule you could point at.

Now pick any single company inside a bid-rigging cartel. Same story. It bids
for public contracts. Sometimes it wins. Sometimes it loses. Perfectly normal.

That's the whole difficulty \bt and it's why this is a hard problem rather than
a big one. Looked at one at a time, there is nothing to find. Nobody looks
guilty on their own.

Because the crime isn't in any one of them. The crime is in the
\emph{arrangement}.

So we stopped looking at them one at a time. We built a system that looks at
who deals with whom \bt and it turns out coordination has a shape you can
actually see.

Our project is called CollusionGraph. One system, screening two completely
different crimes. Dirty money on one side, rigged public contracts on the
other.

Over the next twenty minutes we'll take you from that one idea \bt to the
working console an investigator would sit in front of. Which we built and
tested ourselves.

I'm \blank, and I'll take you through the product.
\end{saythis}

\begin{meaning}
This is the most rehearsed forty seconds of the talk. Everything after it is
easier if this lands, so learn it properly rather than reading it.

\textbf{Why it opens with a puzzle rather than a summary.} A room's attention is
highest in the first fifteen seconds and you only get to spend it once.
``Our project is about collusion detection'' spends it on nothing. ``Nobody
looks guilty on their own'' makes the room want the answer, and then you give
it to them.

\textbf{The line to land, and how.} \emph{``The crime isn't in any one of them.
The crime is in the arrangement.''} Pause before it. Slow down through it. That
single sentence is the thesis of the whole project, and if the examiner
remembers one thing from twenty minutes, this is the thing worth remembering.

\textbf{The shape of the opening,} so you can rebuild it if you lose your
place: a puzzle, the same puzzle again in the other world, the difficulty named
out loud, the turn, and then the promise of where we're going. Puzzle, puzzle,
problem, turn, promise.

\textbf{Two words to be clear about from the start.} \emph{Screening} means
narrowing a huge list down to the few worth a human look --- it is never
deciding somebody is guilty. And the \emph{two sides} are money laundering and
bid rigging; we'll call them ``money'' and ``contracts'' all the way through.

\textbf{Do not oversell.} ``Which we built and tested ourselves'' is a quiet
line, and it is doing more work than any adjective would. Say it flatly.
\end{meaning}

\begin{asked}
\q{Why start with the dashboard instead of your method?} Because the dashboard
is the only part that proves the project is finished. Anyone can draw a method
slide. A working console running on real public data can't be faked --- and
everything in the method is visible inside it, so we get to show it rather than
claim it.

\q{Is this a real product or a prototype?} A research prototype that runs end to
end. Real public datasets, real trained models, real stored results. What it
isn't is deployed at a bank, and we're careful about that difference all the way
through.

\q{Two crimes in one project --- isn't that just doing two things at once?} It
would be, if we'd built two systems. We built one, and the fact that one system
handles both is the actual finding. If the arrangement is what gives collusion
away, then the industry it happens in shouldn't matter much --- and we tested
whether that's true rather than assuming it.
\end{asked}

% ---------------------------------------------------------------------------
\slide{Slide 2 --- The Problem --- Why Two Crimes Belong Together}{HIGH (8/10)}{$\approx$ 1 minute}

\begin{saythis}
So let me show you what that arrangement actually looks like, in each of our
two worlds.

Money laundering works by pushing dirty money through a chain of accounts,
until nobody can see where it started. That needs a chain. One account can't
launder anything on its own \bt there's nowhere for the money to go.

Bid rigging works by companies quietly agreeing in advance who's going to win
which contract, so nobody has to compete and the price stays high. That needs
at least two of them. One company can't rig an auction against itself.

Notice what both of those have in common. The thing that makes it a crime is
the \emph{agreement between them}. And an agreement isn't written down anywhere
in the data. What \emph{is} in the data is the trace it leaves.

Which is why we didn't build this as a spreadsheet problem, where you go down a
list of customers and score each one. We built it as a network problem, where
the thing you're looking at is who deals with whom.

And that one decision is what let us build a single system instead of two.
\end{saythis}

\begin{meaning}
This is the intellectual core of the project and it takes ten seconds to say
badly and one minute to say well. Take the minute.

\textbf{The everyday version.} Think about how you'd spot people cheating in an
exam. You wouldn't look at one student's paper. You'd look at which papers have
the same wrong answers. The cheating is in the \emph{relationship} between
papers, not in any single paper.

That is exactly what we're doing, and it is why an ordinary machine-learning
model --- which looks at one row at a time --- is the wrong tool here.

\textbf{The word to introduce.} A \textbf{network} (some people say graph) is
just dots and lines. Dots are accounts or companies. Lines are payments or
shared contracts. Every dataset we use gets turned into dots and lines, and
after that the system doesn't care which crime it came from.
\end{meaning}

\begin{asked}
\q{Isn't a network just a fancy word for a database table?} No, and the
difference is the point. A table answers ``tell me about this customer.'' A
network answers ``show me everyone this customer is connected to, and everyone
\emph{they} are connected to.'' The second question is the one that finds a
ring.

\q{Couldn't you just look for known bad patterns with rules?} We do, and we
kept the rules --- they're one of the four things scoring in the background.
But rules only find what someone already thought of. The learning part is there
to find shapes nobody wrote a rule for.
\end{asked}

% ---------------------------------------------------------------------------
\slide{Slide 3 --- What You Are About To See}{Medium (5/10)}{$\approx$ 30 seconds}

\begin{saythis}
So this is the console. It's what an investigator would open on a Monday
morning.

Everything on it is real. Real public datasets, real trained models, real
results --- nothing here is a mock-up and nothing is generated live to look
good.

Six tabs across the top, and we'll go through every one of them.

And before I go any further \bt look at the line along the bottom of the
screen. It says this is a screening signal, not a determination of guilt. That
line is on every screen, and it's on every piece of data the system hands out.
I'll come back to why that matters at the end.
\end{saythis}

\begin{meaning}
Short slide, three jobs. You are establishing that the data is real, giving the
room a map of the six tabs, and planting the ethics line early so that when you
close on it in twenty minutes it lands as a design principle rather than a
disclaimer.

Do not linger here. The room wants to see something happen.
\end{meaning}

\begin{asked}
\q{Where does the data come from?} Six public datasets. Bitcoin payment
records, a European public-contracts database, an international auctions
dataset with confirmed cartels, a synthetic banking dataset, and a full
national procurement archive. All open, all citable, none of it private.
\end{asked}

% ===========================================================================
\part{Part 2 --- The frame around every screen (1 minute 45)}

\slide{Slide 4 --- The Domain Switch --- Two Worlds, One Machine}{Medium (6/10)}{$\approx$ 45 seconds}

\begin{saythis}
Top right. This switch is the whole two-crimes idea in one control.

On the left, Financial, and underneath it says A-M-L \bt anti money laundering.
On the right, Procurement, and underneath, bid rigging. We name the actual
crime, because ``financial'' on its own tells you nothing.

\do{click the switch} Watch what happens. The colours change. The datasets
change. The cases change.

Now watch what \emph{doesn't} change. Nothing underneath. Same reader of the
data, same scoring, same grouping, same evidence rules, same screens.

That's the honest test of whether this is one system or two systems in a
trench coat. It's one.
\end{saythis}

\begin{meaning}
This is a small control carrying a large claim, so it earns its own slide.

If the code had two separate pipelines behind that switch, the project would
just be two projects sharing a menu bar. It doesn't. Everything from the point
where data is loaded onwards is shared, and each domain only supplies its own
trained numbers.

\textbf{Why we added the crime names.} Until late in the project the pills said
only ``Financial'' and ``Procurement''. We then checked the whole console and
found the letters A-M-L appeared nowhere except two buried lines on the About
page. A visitor could not tell what the financial side was screening
\emph{for}. Naming it cost four characters. Mention this if you want to show
you audit your own interface.
\end{meaning}

\begin{asked}
\q{Same model for both?} Same architecture, same pipeline, but each domain
trains its own numbers on its own data. What is genuinely shared is the data
format, the splitting rules, the grouping, the evidence layer and the way we
measure. That sharing is what made it possible to test whether one crime's
knowledge transfers to the other --- which we did, and I can show you the
result.
\end{asked}

% ---------------------------------------------------------------------------
\slide{Slide 5 --- The Dataset Picker}{Medium (4/10)}{$\approx$ 35 seconds}

\begin{saythis}
Next to it, the dataset. Two on each side.

On the money side: Bitcoin payments, and then the same Bitcoin world seen as
wallets rather than individual payments. That second one matters --- it's the
same crime at a different zoom level, and the system has to survive that.

On the contracts side: European public contracts, and an international auctions
dataset where competition authorities have already confirmed which companies
were cartels. That one is our hardest test, because somebody official already
knows the answer.
\end{saythis}

\begin{meaning}
Two datasets per side is a deliberate defence. One dataset per domain invites
the reasonable suspicion that your result is a property of that one dataset.

The wallets view is the strongest thing on this slide. Same underlying world,
different definition of ``a thing'', and the pipeline needs no changes to
handle it.
\end{meaning}

\begin{asked}
\q{You said six datasets, but I only see four.} Correct, and it's deliberate.
Two of them have no usable answer key at the level we show cases, so a
confirmation rate calculated from them would be meaningless. We use those two
in the planted-cartel experiments instead, which need no answer key. Showing
them here would imply an accuracy we cannot measure.
\end{asked}

% ---------------------------------------------------------------------------
\slide{Slide 6 --- The Six Tabs and the Line at the Bottom}{Medium (5/10)}{$\approx$ 25 seconds}

\begin{saythis}
Six tabs, and each one keeps its own colour and icon the whole way through.
After a minute of using it, you find a screen by its colour before you read its
name.

One colour is deliberately missing from that set. Coral \bt this red-orange
\bt is reserved. It means ``flagged'' and it means nothing else anywhere in
the product. So when you see coral later on, you already know what it is.
\end{saythis}

\begin{meaning}
Small point, but it shows deliberate interface design rather than decoration,
and it pays off twice later --- on the constellation and on the network
picture, where you never have to explain what coral means because you already
did.
\end{meaning}

% ===========================================================================
\part{Part 3 --- Overview: the opening screen (2 minutes)}

\slide{Slide 7 --- The Four Numbers That Matter}{HIGH (8/10)}{$\approx$ 1 minute}

\begin{saythis}
First screen. Four numbers across the top, and they're deliberately the four a
manager asks before anything else.

How many suspicious groups came out of this dataset. How many of those are in
the top risk band. How many cases your team can actually review. And how many
datasets are loaded.

Now look at the third one, because everybody skips it and it's the most
important number on the screen.

This system does not hand you a pile of work and wish you luck. \emph{You} tell
it how many cases your team can get through this week. And then every accuracy
number in the whole project is measured at exactly that workload.

Because here's the thing \bt any screening system can catch everything, if
you're willing to review everything. An accuracy number with no workload
attached to it is meaningless.
\end{saythis}

\begin{meaning}
This is where you introduce the project's central measuring idea, and you do it
without naming it. The technical name is \emph{Precision@k}, and the ``k'' is
that third tile. Do not say ``Precision at k'' out loud unless asked.

\textbf{The everyday version.} A metal detector that beeps constantly finds
every coin on the beach. It is also useless. The only fair question is: of the
fifty places it told me to dig, how many had a coin?

\textbf{Why it is a human input.} The right number depends on how many analysts
you have this week. It is a staffing decision, not a modelling decision. What
the system owes you is an honest measurement at whatever number you pick ---
which is what the next screen delivers.
\end{meaning}

\begin{asked}
\q{Why not just report accuracy?} Because on this data you can be
ninety-eight percent accurate by saying ``innocent'' to everything. Only about
one in fifty of these is a real case. Accuracy rewards a model for ignoring
exactly the thing we built it to find.

\q{Who decides the review number in real life?} The team's manager, from
headcount. That's why it's a slider and not a setting buried in a config file.
\end{asked}

% ---------------------------------------------------------------------------
\slide{Slide 8 --- The Constellation}{Medium (5/10)}{$\approx$ 40 seconds}

\begin{saythis}
The big panel is every case in one picture. One dot per group. Bigger dot,
bigger group. Colour is the risk band. And you can click straight into any of
them.

Let me be straight with you about what's real here. The dots are real --- every
one is an actual group at its actual rank with its actual score. The
\emph{arrangement} of them is decorative. Position doesn't mean anything.

What I want you to notice is the shape of the crowd. A handful of hot ones at
the top, and then a long quiet tail. That's what a real screening problem looks
like, and it's exactly why that review number matters so much.
\end{saythis}

\begin{meaning}
Volunteering that the layout is decorative is worth more than the panel is. An
evaluator who works it out for themselves starts wondering what \emph{else} is
decorative, and you have lost the room's trust for the rest of the talk.

Say it, move on, and you have bought credibility cheaply.
\end{meaning}

\begin{asked}
\q{So it's just decoration?} The layout is. The data is not, and you can click
any dot and land on that exact case. If you want the true ordering, the next
tab is the true ordering with no styling at all.
\end{asked}

% ---------------------------------------------------------------------------
\slide{Slide 9 --- The Leaderboard and the Move to the Queue}{Medium (3/10)}{$\approx$ 20 seconds}

\begin{saythis}
On the right, the same thing as a plain list. Rank, risk, group size, and the
name of the pattern where we could put a name to it.

If you had ten seconds and one screen, it'd be this one. But let's go into the
screen an investigator actually works in.
\end{saythis}

\begin{meaning}
Pure transition. Keep it short and click through on the last word --- a talk
that moves feels rehearsed even when it isn't.
\end{meaning}

% ===========================================================================
\part{Part 4 --- The Alert Queue: the working screen (4 minutes 30)}

\slide{Slide 10 --- What One Row Actually Is}{VERY HIGH (9/10)}{$\approx$ 55 seconds}

\begin{saythis}
This is the screen somebody would live in. It's a worklist, worst first.

Before we read a row, two things you need to know.

First \bt a row is not a person and it's not a payment. It's a \emph{group}. A
cluster of accounts that deal with each other far more than they deal with
anybody outside.

And that's not a presentation choice, it's forced by the crime. One account
can't launder money on its own. So the thing you investigate has to be the
group, not the individual.

We score every individual account first, and then we roll them up into groups.
Sixty-seven thousand accounts on this dataset become two hundred and fifty-four
cases. That's the difference between a research output and a morning's work.

Second \bt near-identical groups get merged, so you never open the same case
twice under two different reference numbers.
\end{saythis}

\begin{meaning}
The single most common misunderstanding in the whole demo is that a row is a
suspect. Kill it in the first fifteen seconds.

\textbf{The everyday version.} Police don't investigate one phone number. They
investigate the group of numbers that keep calling each other. Same idea.

\textbf{Two words you may need if pushed.} The grouping method is called
\emph{Leiden}; it finds clusters that are much more connected inside than
outside. The merging step removes overlapping groups. Neither name is worth
saying unless somebody asks.

\textbf{The detail that impresses.} The grouping never looks at the answers.
It runs on structure alone --- which is why the identical step also runs on the
dataset where no answers exist at all.
\end{meaning}

\begin{asked}
\q{Why not just show me the top scoring accounts?} Because you'd hand an
analyst a list of individuals and leave them to rebuild the ring by hand, which
is the work we're trying to save. And the same ring would appear five times as
five separate rows.

\q{How do you decide where a group stops?} The clustering does, from the
network's own structure. We also cap group size, because a ``group'' of ten
thousand accounts is not a case anybody can work.
\end{asked}

% ---------------------------------------------------------------------------
\slide{Slide 11 --- The Review Slider --- The Most Honest Control Here}{VERY HIGH (10/10)}{$\approx$ 1 minute 45}

\begin{saythis}
This slider is the most honest thing in the system, so I want to spend a proper
moment on it.

It asks one question. How many cases can your team actually get through? Say
it's fifty. \do{drag to 50}

Now watch this number just to the right of it.

That's telling me that out of those fifty cases, about a third turn out to be
real. Sixteen out of the fifty, to be exact. And that's not a promise or a
forecast \bt it's measured, on cases where somebody already knew the answer.

Now watch what happens when I get greedy. \do{drag to 100} A hundred cases, and
it drops to about a quarter. \do{drag to 200} Two hundred, and it's down to
roughly one in seven.

That is not the system getting worse. That's what happens to every screening
system ever built when you dig further down the list. The good stuff is at the
top.

Being able to show you exactly how fast it falls off \bt that's the point.
\end{saythis}

\begin{meaning}
If the room remembers one slide, make it this one. It is the project's central
measurement made physical.

\textbf{The everyday version.} Picking the first fifty apples off a tree that's
been sorted best-first. The first fifty are good. The next hundred are worse.
Nobody is surprised by that, and nobody should be surprised here.

\textbf{Why the number jumps rather than sliding smoothly.} Because it refuses
to invent anything. It only ever shows workloads the evaluation actually
measured --- fifty, a hundred, two hundred. If you drag to seventy it shows you
the fifty figure, because that is the last one that was really tested.

\textbf{Do not skip the honesty line.} ``Measured, not promised'' is the phrase
that separates this from a sales demo.
\end{meaning}

\begin{asked}
\q{A third seems low.} Two answers and both matter. First, blind guessing on
this data gets about six in a hundred, so a third is roughly five times better
than chance at the same workload. Second, and this is the bigger one \bt about
three quarters of this dataset has \emph{no recorded answer either way}. A case
we couldn't confirm is not the same as a case we got wrong. So every precision
number we publish is a floor, not a ceiling.

\q{Could you tune it to look better?} We could quote the top ten instead of the
top fifty and the number would jump. That's exactly the game this slider exists
to stop. The workload is stated up front and the number is read at that
workload.
\end{asked}

% ---------------------------------------------------------------------------
\slide{Slide 12 --- Show Me One It Did NOT Flag}{VERY HIGH (9/10)}{$\approx$ 1 minute}

\begin{saythis}
Here's the question I'd ask if I were sitting where you are. ``Fine \bt you've
shown me the suspicious ones. Now show me one it did \emph{not} flag, so I can
see the difference.''

These four buttons do exactly that. All, High risk, Medium, and Low or normal.
And each one carries its own count, so you can see the shape of the queue
before you click anything.

Look at this. On the Bitcoin payments data, out of two hundred and fifty-four
groups, nearly a hundred and fifty are rated ordinary. \do{click Low} There
they are. Same scoring, same everything. The system simply doesn't think
they're worth your morning.

And I'll be straight about one thing. \do{switch to contracts} If I flip to
the European contracts data, this Low button is empty and greyed out.

That is not the system failing to recognise a normal company. That dataset is a
research sample where roughly four in every ten entries are cartel cases by
design. There genuinely aren't many ordinary firms in it to show you.
\end{saythis}

\begin{meaning}
This is the highest-value defensive slide in the deck, because it answers the
sceptic's question \emph{before} they ask it. Answering it unprompted is worth
far more than answering it under pressure.

\textbf{Background you need.} Until late in the project the console could not
do this at all. The queue only ever drew the top cases by score, so every row
on screen was a high-risk row, and it looked as though the system could not
recognise normal activity. The ordinary groups were always in the data ---
nothing surfaced them.

\textbf{The empty band.} Volunteer it. A greyed-out button with a zero on it
looks like a bug for about four seconds, and if the examiner spots it before
you mention it, you spend the next two minutes on the back foot.
\end{meaning}

\begin{asked}
\q{Isn't an empty band a bug?} It was the first thing we checked. It follows
from how that dataset was built --- it's a matched study sample, not a real
population. For the contrast case, look at the Bitcoin queue, where more than
half the groups are ordinary.

\q{Where do the band boundaries come from?} They're fixed thresholds on the
score, the same on every dataset, so the bands mean the same thing everywhere.
\end{asked}

% ---------------------------------------------------------------------------
\slide{Slide 13 --- Reading One Row}{HIGH (7/10)}{$\approx$ 1 minute 10}

\begin{saythis}
Left to right across a row.

Rank. Then the risk score \bt and that's a genuine probability, not an
arbitrary number between zero and one. Every model's output is put onto a
proper scale before we compare or combine anything, and I'll show you what
happens if you skip that step.

Then the pattern, if we could name one. There are nine named patterns and they
aren't ours \bt the money ones come from F-A-T-F, the contract ones from the
O-E-C-D. Those are the international bodies that publish what these crimes
actually look like. We only put a name on it when we can prove the shape.
Having no name is an honest answer, not a blank.

Then how many official warning signs matched. Then the size of the group.

Then this little chart, which is activity over time. Flat means everything
happened in one window \bt very common in the Bitcoin data. Wiggly means the
activity is genuinely spread out.

Then when it happened, and the case reference.
\end{saythis}

\begin{meaning}
\textbf{The word to be ready for is ``calibrated''.} What it means: two models
can both be good at ranking and still be on totally different scales. One says
nought point nine when it's fairly confident; the other says nought point one
when it's certain. Average those raw and the loud one flattens the quiet one.

We measured that. Put them on a proper scale first and the combined result is
about nine times better than combining them raw. Same models, same data. That
one ordering decision was worth more than any model we trained.

\textbf{Say the standards bodies out loud.} It moves the nine patterns from
``something a student invented'' to ``something a regulator published'', and
that is a large shift in how the rest of the evidence is received.
\end{meaning}

\begin{asked}
\q{What makes a probability ``genuine''?} If the system says nought point eight
on a hundred cases, roughly eighty of them should be real. We check that on
held-out data. Without that check a score is only a ranking, and you can't
compare two rankings or average them.

\q{Where does the sparkline come from?} From the real timestamps on that
group's connections. When they don't vary it draws a flat line rather than
faking a spread.
\end{asked}

% ---------------------------------------------------------------------------
\slide{Slide 14 --- The Shortcuts on a Row}{Medium (3/10)}{$\approx$ 20 seconds}

\begin{saythis}
Hover any row and three shortcuts appear. Show me the network. Show me the case
file. Or ask the assistant about this one specifically.

That last one carries the case across with it, so nobody's retyping a reference
number into a chat box.
\end{saythis}

\begin{meaning}
Trivial slide, but it sets up the Copilot at the end. When you open the
assistant later, it will already be holding a case, and that looks designed
rather than bolted on --- because it is.
\end{meaning}

% ===========================================================================
\part{Part 5 --- Graph Explorer: seeing the shape (1 minute 45)}

\slide{Slide 15 --- The Picture --- This Is The Whole Argument}{VERY HIGH (9/10)}{$\approx$ 55 seconds}

\begin{saythis}
Now the picture. This is that group, drawn out.

The coral dots are the flagged members --- and I told you earlier coral means
one thing only. The grey dots around them are context: neighbours we drew in so
the group sits in a world instead of floating in space.

Bottom left tells you how much you're looking at. And if a group is too big to
draw honestly, it says so on screen rather than quietly showing you a slice and
letting you assume that's everything.

This is the moment the project usually clicks for people. You can \emph{see}
the coordination. That's the whole argument in one picture \bt coordination has
a shape, and here is the shape.
\end{saythis}

\begin{meaning}
Slow down and let the room look. Silence for two seconds here is worth more
than another sentence.

\textbf{Why the grey context dots exist.} Without them every flagged group
looks equally suspicious, because you can't see what it's \emph{unusually}
connected to. The contrast is the information.

\textbf{Why the network is trimmed.} The server cuts it down to this group plus
one step outwards, inside this group's own time range, before it's sent. A
browser can't draw a million connections. And a silently cropped picture is
worse than a small one, because the viewer assumes it's complete.
\end{meaning}

\begin{asked}
\q{Is this the whole network?} No, and it says so when it isn't. It's this
group's members plus one step of neighbours, cut to this group's time range on
the server.

\q{Could you be cherry-picking a pretty case?} Fair challenge --- pick a row
yourself and I'll open it. Every case in the queue draws the same way.
\end{asked}

% ---------------------------------------------------------------------------
\slide{Slide 16 --- The Replay, and a Bug We Chose Not to Hide}{HIGH (7/10)}{$\approx$ 50 seconds}

\begin{saythis}
And you can replay it. \do{press play} The connections light up in order, so
you watch the case assemble itself.

Now read the label on that button, because it's doing something careful.

Sometimes it says replay \emph{flow}, and that means the connections carry
genuinely different dates and you're watching real time pass. Sometimes it says
replay \emph{order}, and that means everything in this group happened in the
same window, so what you're watching is sequence, not time.

We found that the hard way. In the Bitcoin data, the top case has a hundred and
one connections that all sit at the exact same moment. The original replay
animated strictly by date \bt so with every date identical, nothing moved and
it looked broken.

We could have spread the dates out to make it animate nicely. That would have
been inventing data. So instead it detects the situation and tells you which
one you're watching.
\end{saythis}

\begin{meaning}
This is a ``we found our own bug and fixed it honestly'' story, and those are
disproportionately persuasive with examiners because they cannot be faked.

The underlying point is a principle you can state: \textbf{when the data can't
support the visual, change the visual, don't change the data.}
\end{meaning}

\begin{asked}
\q{Why not interpolate the dates so it always animates?} Because the animation
would look better and mean less. It would be showing a sequence of events that
never happened in that order.
\end{asked}

% ===========================================================================
\part{Part 6 --- Case Detail: where a person decides (4 minutes 45)}

\slide{Slide 17 --- The Case File Header}{Medium (5/10)}{$\approx$ 35 seconds}

\begin{saythis}
This is the case file. It's where somebody actually decides.

The top line reads like a sentence on purpose. Which case, which dataset, and
which model scored it. Underneath in small grey type is the machine reference.
That's what the software and the export and the assistant all use, so it's
there when you need it and out of the way when you don't.

And the score. More than nine tenths on this one, which is about as high as
this system goes.
\end{saythis}

\begin{meaning}
Small but worth having. Machine identifiers used to be the single most
confusing thing on this screen --- a reviewer would ask ``what is
rgcn\_actor\_s0?'' and the answer took a minute. Now the plain sentence is the
headline and the raw identifier is a subtitle.
\end{meaning}

% ---------------------------------------------------------------------------
\slide{Slide 18 --- The Pattern --- Including When There Isn't One}{VERY HIGH (10/10)}{$\approx$ 1 minute 15}

\begin{saythis}
First panel: the pattern we detected, drawn as a diagram.

Now, this particular case is the interesting one, and I'm choosing it
deliberately even though it makes life harder for me. There's no named pattern
here. It doesn't match any of the nine textbook shapes.

That used to leave the screen looking empty on exactly the cases you most need
to justify. A big score, and two blanks underneath it.

So now we write out what actually happened, in plain words. Read it with me.

Eighty-five accounts. The money went through all of them in one go \bt one
account to the next, then the next, then the next. It never came back to where
it started. And all of it happened in one short stretch of time.

That's a describable situation. It doesn't need a textbook name to be worth
twenty minutes of somebody's morning.
\end{saythis}

\begin{meaning}
This slide answers the hardest fair criticism of any machine-learning system:
\emph{it flagged something and can't tell me why}.

\textbf{The distinction to hold on to.} The model flagged it. The
\emph{matcher} --- a separate set of hand-written rules --- couldn't name it.
Those are different things, and the matcher deliberately never guesses; it
only names a shape when it can formally prove it.

\textbf{Why ``eighty-five accounts, eighty-four connections'' means a chain.}
If you have eighty-five things and only eighty-four links between them, there
is no way to form a loop. Money can only travel forwards. That is what licenses
the word ``chain'', and it is measured, not asserted.

\textbf{Language note.} This wording is the third attempt. The first two were
rejected as too technical. The rule now is that these sentences describe the
\emph{situation}, never the system --- so there is no ``the model found'' or
``we asked the computer'' anywhere in them.
\end{meaning}

\begin{asked}
\q{So the model flagged something it can't explain?} It flagged something the
rule-matcher can't \emph{name}. The explanation is right there in plain
sentences, and it's built from measured facts about the group. What we won't do
is invent a pattern name to fill the gap.

\q{How many of your cases have no name?} A lot of them, and that's expected.
Real groups often don't contain a textbook shape. On the contracts data most
top cases are two-company groups, which are too small to contain any shape at
all.
\end{asked}

% ---------------------------------------------------------------------------
\slide{Slide 19 --- The Official Warning Signs}{HIGH (8/10)}{$\approx$ 50 seconds}

\begin{saythis}
Second panel: the official warning signs.

When a case matches one, you don't get our opinion. You get the actual
published indicator, quoted, with the body it came from stamped on it, and
underneath it, why it matched \emph{here}.

And this is the panel I'd defend hardest, so let me say why.

An investigator can throw out everything our neural network says \bt distrust
it completely \bt and this evidence still stands. Because it's a rule anybody
can check by hand against a document a regulator published.

On this particular case there are none. And rather than show you an empty box,
we tell you plainly what did make it stand out.
\end{saythis}

\begin{meaning}
The phrase that matters is \textbf{``even if you don't trust the neural
network''}. It converts the evidence panel from a feature into a safeguard.

Mechanically: these come from hand-written rules over the group's network, not
from the learned model, and the indicator wording is quoted from a fixed
reference file. The learned model has no way to touch this panel.
\end{meaning}

\begin{asked}
\q{Could the model have invented a warning sign?} It has no route to. The
matching is deterministic hand-written rules, and the text is quoted from a
fixed file. Different mechanism entirely.

\q{Who are FATF and OECD?} The Financial Action Task Force sets the
international anti-money-laundering standards. The OECD publishes the standard
checklist for detecting bid rigging in public procurement. Both are the
reference documents actual investigators are trained on.
\end{asked}

% ---------------------------------------------------------------------------
\slide{Slide 20 --- The Evidence Panel}{Medium (6/10)}{$\approx$ 40 seconds}

\begin{saythis}
Third panel: the facts themselves. How many are in the group. How many
connections between them. When it happened. Amounts and fees where the dataset
actually records them.

And every fact is tagged with where it came from.

That matters more than it sounds. If a fact is missing here, it's genuinely
missing. We don't fill in a guess and we don't quietly average something in. On
a dataset that has no amounts, the amounts row simply isn't there.
\end{saythis}

\begin{meaning}
``We never fill in a guess'' is a claim about data handling that examiners
like, because the opposite --- silently imputing missing values --- is one of
the most common quiet errors in student projects.

It's also true here as a design rule: fields a dataset doesn't have stay absent
rather than being filled in with an average.
\end{meaning}

\begin{asked}
\q{Why not impute the missing values?} Because on a screening system an
imputed amount becomes a reason somebody gets investigated. We'd rather show
less and have all of it be real.
\end{asked}

% ---------------------------------------------------------------------------
\slide{Slide 21 --- How Solid Is This? --- Marking Its Own Homework}{VERY HIGH (9/10)}{$\approx$ 1 minute 10}

\begin{saythis}
Fourth panel, and this is my favourite thing in the product, because it's the
system checking its own work in public.

First I need to tell you what it's checking.

When the system raises a case, it also picks out a small handful of connections
out of the whole group. It says: these specific few are what made me raise it.

On this case, that's three accounts and four payments between them. Everything
else in the group is context.

So then there are two obvious questions to ask about that handful, and the
panel asks both in plain English.

One. Is that handful enough on its own? We re-run the case showing the system
\emph{only} those three accounts and four payments, and nothing else. If the
warning still comes out, then yes \bt that really was enough.

Two. Was it actually necessary? We re-run the case again with those same four
payments \emph{deleted}. If the warning disappears, then yes \bt they really
were load-bearing.

Yes or no. For a single case there's no in-between, and dressing that up with a
decimal point would be false precision.

And when the two answers contradict each other, the panel says so, rather than
hiding it.
\end{saythis}

\begin{meaning}
\textbf{Never say ``the highlighted links'' without first saying what they
are.} That was a real complaint about our own report: the phrase was used
before it was ever defined, so the whole explanation collapsed. The first
paragraph of this script exists purely to define it, and you must not cut it.

The plain definition to hold: \textbf{the handful of connections the system
says its decision depended on.} Out of eighty-five accounts, three of them and
four payments.

\textbf{The everyday version.} Someone claims they failed an exam because of
two questions. Test one: show them only those two questions --- do they still
fail? Test two: delete those two questions --- do they now pass? If both
answers are yes, those two questions really were the reason.

\textbf{The technical names, only if pushed.} These two checks are called
\emph{sufficiency} and \emph{necessity}, and in the tool's own output they
appear as fidelity-minus and fidelity-plus. Confusingly, a zero means the best
possible result on one and the worst on the other, which is exactly why the
panel asks the question instead of printing the number.
\end{meaning}

\begin{asked}
\q{How often does that check fail?} It used to fail on thirty-eight of our
fifty published Bitcoin explanations. We replaced the explanation method and
now it fails on one. We kept the old number written down rather than quietly
improving it.

\q{What if it fails on the case I'm looking at?} The panel puts a warning on
it and tells you to treat the attribution cautiously. It doesn't hide it, and
the case still shows all its other evidence.
\end{asked}

% ---------------------------------------------------------------------------
\slide{Slide 22 --- Nothing Is Hidden}{Medium (3/10)}{$\approx$ 25 seconds}

\begin{saythis}
At the bottom, one collapsed section holding every raw value behind this case,
for whoever wants to argue with it. And an export button that hands you the
whole file.

The rule we set ourselves is that the tidy version can never contain less than
the raw version. If a case carries a field we haven't designed a panel for, it
still shows up down here.
\end{saythis}

\begin{meaning}
This is a structural guarantee, not a promise --- unrecognised fields fall
through to that section automatically, so a new field cannot be silently
swallowed by a future change.
\end{meaning}

% ---------------------------------------------------------------------------
\slide{Slide 23 --- The Case With No File At All}{HIGH (7/10)}{$\approx$ 45 seconds}

\begin{saythis}
Let me deliberately open a quiet one from further down the list. \do{open a low
band case} Watch what you get.

Not an error. It says: this group looks ordinary, the system didn't think it
was worth anyone's time, so it didn't write up a file. And it says plainly that
this is the right outcome.

You still get the score, the size, the time window, and you can still go and
look at the network yourself.

The reason full write-ups stop partway down the list is simply that producing
them costs real computing time, and spending it on cases nobody will read is
waste.
\end{saythis}

\begin{meaning}
Another self-caught bug worth telling. The moment the Low band was added, these
cases became reachable for the first time --- and clicking one showed a red
error box saying ``no explanation bundle''. The system was working exactly as
designed and reporting itself as broken.

The lesson, if you want to say it: \textbf{a designed outcome must never be
displayed as a failure.}
\end{meaning}

\begin{asked}
\q{So most of your cases have no explanation?} Most of the ones nobody would
review, yes. Explanations are produced for the head of the queue, which is the
part a team actually works. Producing them for all of them would cost hours of
computing for pages nobody opens.
\end{asked}

% ===========================================================================
\part{Part 7 --- Model Lab: the evidence (4 minutes 15)}

\slide{Slide 24 --- Never A Score On Its Own}{VERY HIGH (10/10)}{$\approx$ 1 minute}

\begin{saythis}
This screen is where the claims live. And every number on it comes out of a
stored results file \bt nothing on this page is worked out in the browser while
you watch.

The big number at the top is the overall score.

Now, on its own, that number is meaningless. And this is the single thing most
people get wrong when they read machine-learning results.

So directly underneath it, we print what \emph{blind guessing} would have
scored on this exact data. And then the ratio between the two.

Here, guessing gets about six in a hundred. So the model isn't just ``point
four seven'' \bt it's several times better than chance on data where the
crime is genuinely rare.

That's the honest form of every result in this project. Never a score on its
own. Always a score next to what doing nothing clever would have got you.
\end{saythis}

\begin{meaning}
\textbf{The everyday version.} Scoring seventy percent on a test means nothing
until you know whether the class average was thirty or ninety-five.

\textbf{Why this matters here specifically.} On the contracts data, guessing
scores about a third, because that sample is unusually full of cartels. On the
Bitcoin data, guessing scores about six in a hundred. So a raw score of nought
point three nine on contracts is a \emph{worse} result than nought point four
seven on Bitcoin, even though the number looks similar. Quoting either alone
would mislead --- and quoting scores alone is the most common way results in
this field get oversold.
\end{meaning}

\begin{asked}
\q{Why is the guessing level so different between datasets?} Because it's just
how common the crime is in that data. Rare crime, low guessing level, and a
much harder problem.

\q{Doesn't that let you pick a flattering dataset?} It would, which is why we
publish both, including the one where our own network model does worse than
guessing.
\end{asked}

% ---------------------------------------------------------------------------
\slide{Slide 25 --- The Workload Table}{HIGH (7/10)}{$\approx$ 45 seconds}

\begin{saythis}
Underneath, the same run broken down by workload. Read one row with me.

If you review this many cases, this share of them turn out to be real. This is
how much of the total wrongdoing you'd have found. And this is your false
alarm rate.

Three rows, three different staffing levels.

This is the table you'd actually take to whoever signs off headcount, because
it turns a research score into a staffing decision.
\end{saythis}

\begin{meaning}
Two words to be able to define if asked. \textbf{Precision} = of the cases you
opened, how many were real. \textbf{Recall} = of all the real cases out there,
how many you found. They pull against each other: review more, find more, but a
smaller share of what you open is real.

The table only lists workloads that were actually measured. It never draws a
row for a number that was not tested.
\end{meaning}

\begin{asked}
\q{Your recall looks low.} It is, at these workloads, and that's the honest
trade. We're optimising for a small team's morning, not for finding everything.
Reviewing more would raise recall and lower precision, and the table shows you
exactly that trade rather than hiding it.
\end{asked}

% ---------------------------------------------------------------------------
\slide{Slide 26 --- The Three Charts, and the Collapse We Didn't Crop}{HIGH (8/10)}{$\approx$ 1 minute 10}

\begin{saythis}
Three charts, and every one of them exports as an image, because these are the
same figures that go into the write-up. No retyping, so the paper and the
screen can't drift apart.

First: how well it scored, period by period, with the guessing line dashed
behind it.

Now look at the collapse after period forty-three. That is real. It's a known
event in this dataset \bt a large marketplace was shut down, and the behaviour
of the whole network changed underneath the model.

We could have cropped the chart at forty-three. Showing it \emph{is} the
finding.

Second: how many of your top picks are real, drawn against workload, with a
marker at wherever you left the slider.

Third: the same thing at group level. And read the note underneath it, because
it governs this whole screen \bt a case we couldn't confirm is not the same as
a case we got wrong.
\end{saythis}

\begin{meaning}
The step-43 collapse is a genuine strength dressed as a weakness. It is why the
model is trained on the past and tested on the future and never the other way
round, and it is why the report warns that our own validation scores are not a
trustworthy guide when the world shifts underneath the model.

If you are asked whether it invalidates the results: it constrains them, and we
state the constraint. A model that had never met a shift would be the more
suspicious result.
\end{meaning}

\begin{asked}
\q{Doesn't that collapse invalidate everything after it?} It constrains it,
and we report per period rather than hiding it in one blended average. It's
also the reason we test strictly forwards in time.

\q{What was the event?} A large darknet marketplace was taken down. The
behaviour it generated stopped, so the patterns the model learned before that
point were partly describing a world that no longer existed.
\end{asked}

% ---------------------------------------------------------------------------
\slide{Slide 27 --- Can You Trust These Numbers?}{VERY HIGH (9/10)}{$\approx$ 1 minute 30}

\begin{saythis}
Last section here, and it's collapsed by default because it's the follow-up
question, not the opening one. Can I trust any of this?

Four answers. \do{expand the section}

One. Every headline model was built five separate times, each with different
starting randomness, and we report the average \emph{and} the spread. If a
result only shows up on one run out of five, we'd rather find that ourselves
than have you find it.

Two. When we say one method beats another, we re-ran that comparison a couple
of thousand times on reshuffled data, and we show you the range the gap landed
in. If that range crosses zero, we don't get to call it a win.

Three, and this is the hardest test in the project. Take one country out
completely. Train without it. Then test on it. Every country takes a turn.

And look at what that shows. On the auctions data it works everywhere. On the
European data it depends on the country \bt and the biggest country actually
fails. We publish that, because an average that hides a failure is a misleading
average.

Four. We deliberately corrupted some of the training answers to see what
breaks. The test answers were never touched.
\end{saythis}

\begin{meaning}
Publishing your own failing fold is the strongest single credibility move in
the talk. Say it plainly and do not apologise for it.

\textbf{Why five runs.} Neural networks start from random numbers. Run one and
you might have got lucky. Five runs with the spread reported tells you whether
the result is real or an accident.

\textbf{Why hold a whole country out.} Otherwise you're only proving it works
on places it has already seen. Holding a country out entirely is the closest
thing to asking ``would this work somewhere new?''
\end{meaning}

\begin{asked}
\q{Which country fails, and why?} The largest one in that pool, and it scores
slightly worse than guessing there. Two likely reasons: it has by far the most
data so it dominates what the model learned elsewhere, and procurement rules
differ by country. We publish every country rather than the average alone.

\q{Isn't corrupting your own training data strange?} It's a robustness test. If
a small amount of wrong answers destroys the model, it would be useless in the
real world, where labels are always partly wrong.
\end{asked}

% ===========================================================================
\part{Part 8 --- About, the assistant, and the close (3 minutes 30)}

\slide{Slide 28 --- The About Screen}{Medium (6/10)}{$\approx$ 50 seconds}

\begin{saythis}
The About screen is the argument in one page, and it's the one I'd leave up if
we ran out of time.

Two ledgers, one shape.

Then the nine patterns, each with its own diagram. Four money ones, four
contract ones, and one that lives in both worlds \bt hidden common ownership,
where two companies that look independent are quietly owned by the same person.

And I want to stress: those nine are not ours. They're taken from the
international standards. Our contribution is a detector that finds them, and
when we plant them in data ourselves, it recovers all nine.

Then four steps: read the data, learn, rank, explain.

And at the bottom, the boundary. This produces screening signals for human
investigators. Inputs to due process \bt never a substitute for it.
\end{saythis}

\begin{meaning}
The ``these nine are not ours'' line is doing real work. It moves the pattern
library from something a student invented to something a regulator published,
and it makes the recall claim meaningful.

The recall claim itself: the matcher and the thing that plants test patterns
are two independently written pieces of code, so when one finds everything the
other planted, they are genuinely checking each other rather than agreeing with
themselves.
\end{meaning}

% ---------------------------------------------------------------------------
\slide{Slide 29 --- The Assistant, and Trying To Break It}{VERY HIGH (9/10)}{$\approx$ 2 minutes}

\begin{saythis}
Last thing, and it's the part people most want to poke at. There's an
assistant, and you can open it on any screen.

Let me ask it something ordinary first. \do{type: which five cases should I look
at first and why?}

Watch the panel while it thinks. That trace is real \bt those are the actual
lookups it's making, live. It cannot answer from memory. It has to go and read
the published results, and you can open every lookup it made and see exactly
what came back.

Now the badges along the top. Every number in that answer had to appear
literally in something it looked up. If one didn't, it tells you. If it talks
about warning signs without consulting the reference documents, it tells you
that too.

Now let me try to break it. \do{type: is firm X guilty?}

There. ``This system does not determine guilt.'' It won't answer that question.
And that isn't a filter on the way in \bt it's checked on the way out, on every
single response.

One more. \do{type: what will the score be in 2027?} It refuses, and it lists
what it actually has instead of inventing something.

Before any release, it sits a twenty-four question exam covering exactly these
failure modes. All twenty-four have to pass.
\end{saythis}

\begin{meaning}
Rehearse this one. It is the only live, unscripted-looking moment in the talk
and it is the most impressive if it works.

\textbf{Have a screenshot ready} in case the key or the network fails. If it
does fail, say so plainly and show the screenshot --- a confident recovery
costs nothing, a panicked one costs the room.

\textbf{The point to land.} The guarantees live \emph{outside} the language
model. Swap the model tomorrow and the rules still hold, because they are
checks applied to its output, not instructions given to it.

\textbf{If there is no key on the machine} it says so plainly rather than
pretending, and every other part of the console still works.
\end{meaning}

\begin{asked}
\q{What stops it making a number up?} A check that runs after it has spoken.
Every number in the answer must appear literally in the evidence it retrieved.
If it doesn't, the answer is marked unverified on screen rather than shown
clean.

\q{Which model is behind it?} A hosted open model. But the important answer is
what it's \emph{allowed} to do: read-only lookups against published results,
nothing else. It cannot change a score and it cannot see anything a user
couldn't see.

\q{Isn't an AI assistant the riskiest thing you could add here?} It would be,
unguarded. That's precisely why it's the most heavily tested component in the
project.
\end{asked}

% ---------------------------------------------------------------------------
\slide{Slide 30 --- The Close}{HIGH (8/10)}{$\approx$ 40 seconds}

\begin{saythis}
One last thing and then I'll stop.

That line along the bottom has been on every screen for the last twenty
minutes. And it's on every piece of data the system hands out, not just the
pretty pages. Screening signals for human investigators. Not determinations of
wrongdoing.

That isn't a disclaimer we added at the end to be safe. It's the reason the
assistant refuses certain questions. It's why the case files show you their own
weaknesses. And it's why we publish the country where it fails right next to
the ones where it works.

A tool that decided who was guilty would be a different project, and a much
worse one. This one shortens the list a person has to read.

Thank you \bt happy to take questions.
\end{saythis}

\begin{meaning}
Close on the ethics boundary as a design driver rather than a legal shield.
Every claim in this block is checkable on screen within ten seconds, which is
what makes it the right note to end on.

Then stop talking. Do not add a summary slide after this.
\end{meaning}

% ===========================================================================
\part{Appendices}

\slide{Appendix A --- The Five Minute Version}{---}{$\approx$ 5 minutes}

If you are cut short, run slides \textbf{1, 11, 12, 15, 18, 29, 30}.

That is: the claim, the honest measurement, the ordinary case, the picture, the
explanation, the guardrails, and the boundary. It still contains a complete
argument.

\slide{Appendix B --- Never Read A Decimal Aloud}{---}{Reference}

\begin{center}\small
\begin{tabular}{@{}p{6.2cm}p{8.4cm}@{}}
\hline
\textbf{On screen} & \textbf{Say} \\
\hline
Precision 0.32 at 50 cases & ``about a third --- sixteen of the fifty'' \\
Precision 0.23 at 100 & ``about a quarter'' \\
Precision 0.135 at 200 & ``roughly one in seven'' \\
Guessing level 0.065 & ``about six in a hundred'' \\
148 of 254 below 0.2 & ``nearly a hundred and fifty of the two hundred and
fifty-four'' \\
Contracts, 2 of top 4 & ``two of the top four'' \\
Auctions, Italy, 3 cases & ``three cases, and all three are confirmed
cartels'' \\
Risk score 0.9265 & ``more than nine tenths'' \\
77\% unlabelled & ``about three quarters has no recorded answer either way'' \\
Transfer lift 1.57 & ``about half again better than guessing'' \\
Failing fold 0.90 & ``actually slightly worse than guessing there'' \\
\hline
\end{tabular}
\end{center}

\slide{Appendix C --- Questions That Have Actually Been Asked}{---}{Reference}

\begin{enumerate}
  \item \textbf{``Show me one it didn't flag.''} Slide 12. Low band, Bitcoin
        payments, nearly a hundred and fifty of them.
  \item \textbf{``Why is precision only a third?''} Slide 11. Guessing gets six
        in a hundred, and three quarters of the data has no recorded answer, so
        every figure we publish is a floor.
  \item \textbf{``Your simple model beats your neural network --- so why the
        deep learning?''} On one dataset, yes, and we publish it. That dataset
        already has network information copied into its columns, so the simple
        model gets the network for free. In the planted-cartel experiments,
        where there is no answer key at all, the simple model cannot run and
        the deep model finds about nine in ten of the bid-together cartels.
        That experiment is the one closest to real deployment.
  \item \textbf{``Could an investigator be misled by this?''} Slides 21 and 30.
        The case file shows its own weaknesses, the assistant refuses guilt
        questions, and the boundary is on every screen.
  \item \textbf{``Is any of this live data?''} No. Published public datasets,
        scored offline, served read-only. The serving side physically cannot
        retrain or change a score --- that's the trust boundary, and it's the
        centre of the architecture talk.
\end{enumerate}

\end{document}
